\ficha{Weller (2015), Rothschilds' Delicate and Difficult Task}{weller2015}

\begin{identificacao}
WELLER, Leonardo. Rothschilds' ``Delicate and Difficult Task'': reputation,
political instability, and the Brazilian rescue loans of the 1890s.
\textit{Enterprise \& Society}, v. 16, n. 2, p. 381-412, 2015.
DOI: 10.1017/eso.2014.32.

\smallskip
Artigo de periódico, 32 páginas, lido na íntegra, incluindo o apêndice de dados
fiscais e de comércio. Idioma: inglês. Versão lida: a publicada.
\end{identificacao}

\bloco{Problema e tese}

Por que a casa Rothschild de Londres socorreu o Brasil nos anos 1890 em vez de
deixá-lo suspender pagamentos, e por que só em 1898 impôs condições efetivas de
política econômica? A resposta de Weller é a dependência recíproca. O banco
precisava do Brasil para preservar a reputação de emissor sem default, porque a
dívida brasileira respondia por 22,9\% de tudo o que a casa tinha subscrito e
ainda circulava em Londres em julho de 1898 (p. 397). O Brasil precisava do
banco porque nenhuma outra casa lhe emprestava. Essa exposição converteu-se em
poder de mercado indireto do devedor: o governo pôde financiar guerras e
descentralizar tributos sem dar calote, pegando carona no guarda-chuva de
crédito do banqueiro. A pressão ortodoxa só veio quando a política interna se
pacificou e a elite dirigente já queria a ortodoxia.

\chave{was not one in which a strong bank unilaterally imposed harsh lending
conditions on a weak state}{p. 407}

\bloco{Argumento}

\begin{enumerate}[nosep]
\item A relação nasce em 1825, sobrevive ao resgate de 1829 e se institucionaliza
em 1854, quando o governo nomeia a casa agente financeiro na Europa, com
exclusividade sobre a emissão de dívida, monopólio que durou cinco décadas
(p. 388). Nos anos 1890 o Brasil é o maior tomador da casa, num momento em que
ela reduzia drasticamente suas emissões (p. 391).

\item A crise da década é política, não comercial. Golpe de 1889, Encilhamento
sob Rui Barbosa, Revolução Federalista, revoltas da marinha e Canudos elevam o
gasto militar, enquanto a Constituição de 1891 transfere aos estados o imposto
de exportação. O serviço mínimo da dívida externa passa de 9,2\% para 21,5\% da
receita tributária entre 1889 e 1898 (p. 394). Weller insiste que o preço do
café estava alto quando a capacidade de pagamento se deteriorou e baixo quando
ela melhorou (p. 394).

\item Dois contrafactuais fiscais sustentam o diagnóstico. Mantido o orçamento
militar de 1890 a 1892, o gasto extra de guerra, 13,4\% da receita tributária,
teria bastado para o serviço da dívida; mantida a tributação federal sobre
exportações, o Tesouro teria arrecadado 7,2 milhões de libras a mais entre 1893
e 1897, o equivalente a 13,8\% da receita anual (pp. 398-399).

\chave{In strict fiscal terms, Brazil could have avoided the debt crisis}{p. 399}

\item Nos empréstimos de 1893 e 1895 o banco aconselhou contração monetária e um
imposto de importação cobrado em ouro, mas não criou mecanismo de verificação.
Rodrigues Alves prometeu resgatar papel-moeda, voltou a emitir em 1896 e 1897 e
alegou que só o Congresso podia cobrar tributo em ouro (p. 400). Weller lê a
frouxidão como cálculo: cobrar demais de um devedor em guerra civil aumentaria o
risco de retaliação, isto é, de moratória unilateral, que destruiria a marca do
banco. É o problema do orçamento flexível de Boot aplicado à dívida soberana.

\item Em 1898 a ameaça de moratória é usada como alavanca pelo próprio devedor.
Prudente de Morais escreve a Bernardino de Campos que convém insinuar a
suspensão de pagamentos para mudar a posição dos banqueiros, e admite um mês
depois que a suspensão seria inevitável sem o empréstimo (p. 402). O resultado é
o funding loan: 8,6 milhões de libras em bônus no lugar do pagamento de juros
até 1901 e dispensa de amortização até 1911, contra condições agora explícitas e
autoexecutáveis, incineração de papel-moeda em montante equivalente aos bônus
emitidos, novo imposto de importação em ouro para financiar o resgate, e
fiscalização da queima pela filial carioca do London and River Plate Bank, banco
britânico e não brasileiro. Descumprida a condição, o empréstimo seria suspenso
(p. 402).

\item A imposição não foi unilateral. O banco pretendia incinerar até a taxa de
16 pence por mil-réis, o que reduziria M2 em 34\%; Bernardino de Campos alegou
embaraços ao Tesouro e obteve 18 pence, com queda de 31\% (p. 404). E a virada
ortodoxa já vinha de dentro: Campos Sales visitou os diretores da casa em
Londres antes de tomar posse e nomeou Joaquim Murtinho, senador que fizera
carreira criticando a emissão dos anos 1890 (p. 404).
\end{enumerate}

\bloco{Conceitos e definições operacionais}

Reputação é a marca sem default do subscritor, e é ela, não a carteira, que
explica o socorro: a dívida brasileira era em média 4,9\% dos ativos do banco
entre 1889 e 1898, e um calote não o quebraria (p. 398). Exposição reputacional
é medida pela fração dos títulos subscritos pela casa e ainda em circulação na
Bolsa de Londres. Prêmio de risco é a diferença entre o retorno até o vencimento
e o retorno da dívida britânica, escolhido por captar diferenças de prazo entre
os empréstimos (p. 387). Capacidade de pagamento é a razão entre serviço exigido
e receita tributária, e não entre serviço e despesa, métrica que Weller recusa
em Ferguson por inflar o denominador com o gasto de guerra (p. 394). Funding
loan não é default, e sim recontratação de dívida, porque não houve suspensão
unilateral e a Corporation of Foreign Bondholders não classificou o Brasil como
inadimplente (p. 403). Corretor de poder, enfim, nomeia o papel do banqueiro que
reforça um dos lados de uma disputa doméstica já existente.

\bloco{Evidência e método}

Estudo de caso narrativo com apoio quantitativo, cobrindo de 1822 a 1902, com
foco em 1889 a 1898. Fontes primárias de arquivo: Rothschild Archive em Londres,
Biblioteca Histórica do Ministério da Fazenda e Instituto Histórico e Geográfico
Brasileiro no Rio de Janeiro, e Archive Historique Paribas em Paris, sobretudo
correspondência entre ministros e banqueiros, em boa parte inédita. Fontes
quantitativas: o Ledger Balance da casa, o Balanço da receita e despesa da
República, Bouças, as séries do IBGE e o Investor's Monthly Manual. Imprensa
britânica de época como termômetro de mercado, Economist, Times, Investor's
Review e Herapath's Railway. As técnicas são o cálculo do prêmio de risco por
empréstimo, a medida da exposição reputacional em cada data e os dois
contrafactuais fiscais descritos acima.

\bloco{Agentes e vertentes}

Do lado do credor: a casa Rothschild de Londres como agente financeiro
monopolista; Barings e Paribas, que recusam crédito à república recém-proclamada
em 1889; o Crédit Lyonnais, que recusa em 1898; a Corporation of Foreign
Bondholders; o London and River Plate Bank como fiscal do acordo; a empresa
Greenwood, interessada no arrendamento da Estrada de Ferro Central. Do lado
brasileiro: Rui Barbosa e o Encilhamento, Floriano Peixoto e a primeira
tentativa de contração, Rodrigues Alves e Bernardino de Campos como
negociadores, Azevedo Castro como agente do Tesouro em Londres, Prudente de
Morais como autor da ameaça de moratória, Campos Sales e Joaquim Murtinho como
face ortodoxa da pacificação. Weller se filia à literatura de relationship
banking, Boot, Flandreau e Flores, e discute Summerhill, Abreu, Franco, Fritsch,
Topik e Triner.

\chave{Rothschilds played the role of a power broker in favor of leaders
committed to orthodox policies}{p. 405}

\bloco{Críticas e limites}

O próprio autor reconhece que não há evidência de que o banco soubesse, em 1898,
que o país se estabilizava, e admite que a casa pode simplesmente ter tido sorte
(p. 407). A assimetria de informação inerente à relação de agência fica nomeada,
não resolvida. O contrafactual do imposto de exportação supõe que mudanças de
tarifa não afetariam o volume exportado, hipótese apoiada na inelasticidade do
café, mas ainda assim uma hipótese (p. 399). A divergência com Ferguson sobre a
gravidade da crise é de métrica, não de arquivo. O artigo encerra a narrativa
nos primeiros anos 1900 e não trata da Caixa de Conversão nem do debate público
brasileiro: a imprensa que ele lê é a de Londres, e as vozes brasileiras que
aparecem são de gabinete, não de jornal.

\begin{dialogo}
\item Procurar nos quatro diários, em 1906 e entre 1910 e 1913, as palavras
incineração, queima de papel-moeda, funding, resgate e Murtinho, e verificar se
a contração de 1898 a 1902 aparece como precedente virtuoso a preservar ou como
sacrifício a não repetir. Testa a tese de Weller de que a ortodoxia de Campos
Sales e Murtinho teve dono doméstico e não foi ordem externa. Se os jornais a
narram como imposição de Londres, o achado é sobre o discurso de época, não
sobre o fato, e a divergência precisa ser dita assim.
\item Procurar como os jornais nomeiam os Rothschild: agentes financeiros,
credores, banqueiros de Londres, ou potência que dita política ao país. Weller
descreve o banqueiro como parte interessada e negociadora, presa ao cliente pela
reputação. A ponte testa se o vocabulário jornalístico de 1906 a 1914 sustenta
essa figura ou a de um poder externo unilateral, e ajuda a distinguir a agência
financeira histórica da agência de Londres proposta para a própria Caixa.
\item Rastrear se a barganha de 1898 sobre a taxa, 16 contra 18 pence, é
invocada no debate de 1906 sobre 12, 15 ou 27 pence. Weller documenta que a taxa
de conversão foi objeto de negociação explícita entre Tesouro e banqueiro, com o
ministro alegando embaraços ao Tesouro (p. 404). Se os jornais retomam o
precedente, a escolha de 15 pence aparece como continuidade de um método de
barganha, e não como invenção do Convênio de Taubaté.
\item Procurar no Correio Paulistano e em O Paiz o argumento de que a
valorização do mil-réis arruína o exportador. Weller registra apreciação de
66\% entre 1898 e 1903, recessão e corridas bancárias em São Paulo (p. 405).
Testa a hipótese da dissertação de que o eixo de 1906 é valorização contra
estabilidade a taxa nova, com memória fresca do custo da apreciação, e não
metalismo contra papelismo.
\item Verificar se a defesa da Caixa se apoia na necessidade de manter o crédito
externo. Weller mostra que a ortodoxia rendeu 59 milhões de libras em novos
títulos nos anos 1900 e crédito ao próprio estado de São Paulo, inclusive para a
estocagem de café iniciada em 1906 (p. 406). A hipótese é que o argumento do
crédito em Londres circule na imprensa como razão a favor da Caixa.
\end{dialogo}

\usonocapitulo{Parte II. Sustenta a afirmação de que o enforcement da ortodoxia
sobre a periferia devedora não se descreve bem como ordem externa unilateral,
porque o banqueiro tinha interesses próprios, estava preso ao cliente pela
reputação e agiu como corretor de poder a favor de uma coalizão doméstica já
formada. Serve também para datar e nomear o repertório que os jornais de 1906 a
1914 herdam: funding loan, incineração, agência financeira, Murtinho, Campos
Sales, taxa de conversão negociada em pence. O artigo não trata da Caixa de
Conversão, de modo que a ponte com as fontes primárias é de mecanismo e de
repertório, não de objeto.}
